\documentclass[12pt,a4paper]{article}
\usepackage{/home/maenwe/Documents/Carnet_a_idee/Idee_en_bazar/Aide_a_la_redaction/Notes/package_perso}

\begin{document}
\noindent\textbf{ Note 0}
\begin{enumerate}[label=$-$]
\item Type : équation
\item Contenu complet ? Oui
\item Contenu :\\
\hspace*{.5cm}\begin{minipage}{20cm}
     $2x=x+1$ 
\end{minipage}
\item Note(s) qui l'implique : -
\item Note(s) qu'elle implique : 1
\end{enumerate}
\vspace{.2cm}
\hrule

\noindent\textbf{ Note 1}
\begin{enumerate}[label=$-$]
\item Type : équation
\item Contenu complet ? Oui
\item Contenu :\\
\hspace*{.5cm}\begin{minipage}{20cm}
     $4x=2x+2$ 
\end{minipage}
\item Note(s) qui l'implique : 0
\item Note(s) qu'elle implique : 2
\end{enumerate}
\vspace{.2cm}
\hrule

\noindent\textbf{ Note 2}
\begin{enumerate}[label=$-$]
\item Type : démonstration
\item Contenu complet ? Non
\item Contenu :\\
\hspace*{.5cm}\begin{minipage}{20cm}
     Ceci implique plus de travail à l'aide de la note 2 et ::1::.
\end{minipage}
\item Note(s) qui l'implique : 1
\item Note(s) qu'elle implique : -
\end{enumerate}
\vspace{.2cm}
\hrule


\end{document}